%======================================================================
% SANDRA ECOSYSTEM - Diseño Revista Editorial (Sin Imágenes)
%======================================================================

\documentclass[11pt, spanish, letterpaper, openany]{book}

\usepackage{fontspec}
\setmainfont{Inter}[Path = assets/fonts/, Extension = .otf, UprightFont = *-Regular, BoldFont = *-Bold]

\usepackage[
    top=2cm,
    bottom=2.5cm,
    left=2cm,
    right=2cm,
    columnsep=16pt
]{geometry}

\usepackage{graphicx}
\usepackage{fancyhdr}
\usepackage{titlesec}
\usepackage{xcolor}
\usepackage{microtype}
\usepackage{tikz}
\usetikzlibrary{calc}
\usepackage[most]{tcolorbox}
\usepackage{parskip}
\usepackage{booktabs}
\usepackage{setspace}

%--------------------------------------------------------------------
% COLORES
%--------------------------------------------------------------------
\definecolor{sandraDark}{HTML}{1A2634}
\definecolor{sandraTeal}{HTML}{00897B}
\definecolor{sandraGold}{HTML}{FF8F00}
\definecolor{sandraLight}{HTML}{E0F2F1}
\definecolor{sandraCream}{HTML}{FAFAFA}
\definecolor{sandraGray}{HTML}{90A4AE}
\definecolor{sandraText}{HTML}{455A64}
\definecolor{isoBlue}{HTML}{1565C0}
\definecolor{isoOrange}{HTML}{EF6C00}
\definecolor{isoAmber}{HTML}{FFF8E1}

%--------------------------------------------------------------------
% ESPACIADO
%--------------------------------------------------------------------
\setstretch{1.5}
\setlength{\parskip}{1ex}

%--------------------------------------------------------------------
% CABECERA
%--------------------------------------------------------------------
\pagestyle{fancy}
\fancyhf{}
\fancyhead[RO]{\rightmark}
\fancyhead[LE]{\leftmark}
\fancyfoot[CO]{\thepage}
\fancyfoot[RO,LE]{\color{sandraGray}\fontsize{7}{9}\selectfont\itshape Sandra Ecosystem}

\addtolength{\headheight}{14pt}

%--------------------------------------------------------------------
% BANNER
%--------------------------------------------------------------------
\newcommand{\chapterbanner}[1]{
\thispagestyle{empty}
\begin{tikzpicture}[remember picture, overlay]
    \fill[sandraDark] (current page.north west) -- (current page.north east) -- (current page.south east) -- (current page.south west) -- cycle;
    \fill[sandraGold] ([xshift=0pt] current page.north west) rectangle ([xshift=5mm] current page.south west);
    \node[anchor=south east] at ([xshift=-1cm, yshift=1cm] current page.south east) {
        {\fontsize{120}{130}\selectfont\bfseries\color{white!10}#1}
    };
    \node[anchor=west] at ([xshift=1.2cm, yshift=-2.5cm] current page.north west) {
        \parbox{0.6\linewidth}{
            \fontsize{32}{38}\selectfont\bfseries\color{white}#1
        }
    };
    \fill[sandraTeal] (current page.north west) rectangle ([xshift=5mm, yshift=-3mm] current page.north west);
\end{tikzpicture}
\clearpage
}

\titleformat{\section}{\normalfont\large\bfseries\color{sandraDark}}{\thesection}{0.5em}{}
\titlespacing*{\section}{0pt}{1.5ex plus 0.5ex}{0.5ex}

%--------------------------------------------------------------------
% CAJAS
%--------------------------------------------------------------------
\newtcolorbox{notaeditorial}{
    colback=sandraCream,
    colframe=sandraTeal,
    boxrule=0.5pt,
    left=6pt,
    right=6pt,
    top=6pt,
    bottom=6pt,
    before skip=8pt,
    after skip=8pt
}

\newtcolorbox{estandarISO}{
    colback=sandraLight,
    colframe=isoBlue,
    boxrule=0.5pt,
    rounded corners,
    before skip=8pt,
    after skip=8pt
}

\newtcolorbox{alertaISO}{
    colback=isoAmber,
    colframe=isoOrange,
    boxrule=0.5pt,
    rounded corners,
    before skip=8pt,
    after skip=8pt
}

%--------------------------------------------------------------------
% DROP CAP
%--------------------------------------------------------------------
\usepackage{lettrine}
\lettrinefont{\fontsize{32}{36}\selectfont\bfseries\color{sandraDark}}
\lettrineafterlen=2pt
\lettrinelinenums=false

%--------------------------------------------------------------------
% SEPARADOR
%--------------------------------------------------------------------
\newcommand{\separador}{
    \vspace{1ex}
    \begin{center}
        \color{sandraTeal}\rule{0.12\linewidth}{0.5pt}
    \end{center}
    \vspace{1ex}
}

\usepackage[spanish]{babel}

\begin{document}

%======================================================================
% CAPÍTULO 2: MARCO TEÓRICO
%======================================================================
\chapterbanner{02}

\section{Estado del Arte}

\lettrine{E}{l} middleware empresarial ha evolucionado significativamente hasta convertirse en plataformas complejas de integración. El mercado ofrece soluciones consolidadas como IBM Integration Bus, MuleSoft, Apache Camel y Spring Integration, cada una con fortalezas específicas para diferentes escenarios de uso.

Sin embargo, la mayoría fueron diseñadas para mercados desarrollados, ignorando las necesidades específicas de regiones con infraestructura limitada y contextos regulatorios distintos. Esta brecha crea una oportunidad clara para soluciones localmente desarrolladas que aborden estos desafíos particulares.

\begin{notaeditorial}
    \textbf{Nota del Autor}: El desarrollo de middleware propio representa una oportunidad estratégica para garantizar la soberanía tecnológica en entornos gubernamentales y empresariales críticos.
\end{notaeditorial}

\separador

\section{Fundamentos de Middleware}

El middleware actúa como puente entre aplicaciones y el sistema operativo, facilitando la comunicación e integración entre sistemas heterogéneos. Un Enterprise Service Bus (ESB) representa la evolución hacia una arquitectura basada en servicios, proporcionando un medio centralizado para la integración.

\begin{estandarISO}
    \textbf{Implementación en Sandra}: Server utiliza Go para proporcionar un ESB de alto rendimiento capaz de manejar miles de conexiones simultáneas con mínimo consumo de recursos.
\end{estandarISO}

Las características principales del middleware moderno incluyen: transformación de mensajes entre diferentes formatos; enrutamiento basado en reglas de negocio; composición de servicios complejos; y manejo de transacciones distribuidas.

\separador

\section{Arquitectura de Microservicios}

La arquitectura de microservicios construye aplicaciones como servicios pequeños, autónomos y débilmente acoplados. Cada microservicio es responsable de una funcionalidad específica y puede desplegarse independientemente, permitiendo escalabilidad granular y tolerancia a fallos.

\begin{estandarISO}
    \textbf{Ecosistema Sandra}: Server (middleware central), Consola (interfaz web), SDC (aplicación de escritorio), y Sentinel (procesamiento especializado de eventos).
\end{estandarISO}

\begin{table}
    \centering
    \small
    \caption{Componentes del ecosistema}
    \begin{tabular}{lp{5cm}}
        \toprule
        \textbf{Componente} & \textbf{Descripción} \\
        \midrule
        Server & Middleware central de comunicaciones \\
        Consola & Interfaz administrativa web \\
        SDC & Aplicación de escritorio multiplataforma \\
        Sentinel & Procesamiento especializado de eventos \\
        \bottomrule
    \end{tabular}
\end{table}

\separador

\section{Modelo Zero Trust}

El modelo Zero Trust elimina la confianza implícita, verificando cada solicitud independientemente de su origen en la red. Este paradigma asume que ningún usuario o dispositivo es confiable por defecto, requiriendo autenticación y autorización continuas para cada acceso.

\begin{estandarISO}
    \textbf{ISO 27001}: Sistema de Gestión de Seguridad de la Información
\end{estandarISO}

Los principios fundamentales del modelo son: verificación explícita de usuarios y dispositivos mediante autenticación multifactor; acceso con mínimo privilegio, otorgando solo los permisos necesarios; y verificación continua de seguridad mediante monitoreo de sesiones y detección de anomalías.

\begin{alertaISO}
    \textbf{Consideración importante}: La implementación de Zero Trust requiere un cambio cultural organizacional, no solo implementación técnica.
\end{alertaISO}

Sandra implementa estos principios en todas las capas: cada solicitud de API es autenticada y autorizada, las comunicaciones usan TLS mutuo, y las sesiones tienen duración limitada.

\separador

\section{Tecnologías del Ecosistema}

El ecosistema utiliza tecnologías de última generación, seleccionadas por su rendimiento, seguridad y comunidad activa:

\begin{itemize}
    \item \textbf{Go 1.24}: Rendimiento y concurrencia nativa
    \item \textbf{Rust 2021}: Seguridad de memoria y rendimiento
    \item \textbf{Angular 16+}: Framework frontend enterprise
    \item \textbf{Tauri 2.0}: Aplicaciones de escritorio ligeras
    \item \textbf{gRPC}: Comunicación entre servicios de alto rendimiento
    \item \textbf{WebSockets}: Comunicación tiempo real bidireccional
\end{itemize}

\begin{notaeditorial}
    \textbf{Nota Técnica}: La combinación de Go para el backend y Tauri para el frontend proporciona un equilibrio óptimo entre rendimiento y consumo de recursos.
\end{notaeditorial}

\clearpage

%======================================================================
% CAPÍTULO 3: MARCO LEGAL
%======================================================================
\chapterbanner{03}

\section{Ley de Infogobierno}

\lettrine{L}{a} Ley de Infogobierno establece normas para el uso de tecnologías de información en el sector público venezolano, promoviendo la Soberanía Tecnológica y el uso de Software Libre. Esta legislación representa un avance significativo hacia la reducción de la dependencia tecnológica extranjera.

\begin{table}
    \centering
    \small
    \caption{Cumplimiento con la Ley de Infogobierno}
    \begin{tabular}{lp{5cm}}
        \toprule
        \textbf{Requisito} & \textbf{Implementación} \\
        \midrule
        Software Libre & Todas las tecnologías son código abierto \\
        Estándares Abiertos & REST, gRPC, JSON, Protocol Buffers \\
        Interoperabilidad & APIs documentadas, adaptadores modulares \\
        Seguridad & Cifrado AES-256, JWT, MFA, TLS 1.3 \\
        \bottomrule
    \end{tabular}
\end{table}

\separador

\section{Ley Especial contra los Delitos Informáticos}

Esta ley establece conductas penalizadas relacionadas con el uso indebido de sistemas informáticos. Sandra implementa controles que ayudan a prevenir y detectar estas conductas:

\begin{itemize}
    \item \textbf{Acceso no autorizado}: Autenticación multifactor (MFA)
    \item \textbf{Interceptación de datos}: Cifrado en tránsito y en reposo
    \item \textbf{Falsificación de datos}: Firmas digitales y logs de auditoría
    \item \textbf{Destrucción de datos}: Sistemas de backup y redundancia
\end{itemize}

\separador

\section{Estándares Internacionales}

\begin{estandarISO}
    \textbf{ISO/IEC 27001}: Sistema de Gestión de Seguridad de la Información (SGSI). Proporciona un enfoque sistemático para gestionar información sensible mediante controles técnicos, organizacionales y legales.
\end{estandarISO}

\begin{estandarISO}
    \textbf{ISO/IEC 27002}: Controles de seguridad de la información. Incluye 93 medidas organizadas en cuatro categorías: organizacionales, personales, físicos y tecnológicos.
\end{estandarISO}

\begin{estandarISO}
    \textbf{NIST SP 800-53}: Controles de seguridad y privacidad del Instituto Nacional de Estándares y Tecnología de Estados Unidos.
\end{estandarISO}

\separador

\section{Licenciamiento Soberano}

La elección de licencias en el Proyecto Sandra no es solo legal, es estratégica. Se priorizan licencias que garantizan la autonomía tecnológica y la protección de la propiedad intelectual en entornos críticos orientados a la defensa y la soberanía del Estado.

\begin{alertaISO}
    \textbf{Nota Estratégica}: Las licencias permisivas como Apache 2.0 y MIT permiten que innovaciones basadas en código abierto puedan ser privatizadas y clasificadas, un requerimiento fundamental para el sector defensa.
\end{alertaISO}

\begin{table}
    \centering
    \small
    \caption{Análisis de licenciamiento}
    \begin{tabular}{lp{5cm}}
        \toprule
        \textbf{Licencia} & \textbf{Implicación} \\
        \midrule
        Apache 2.0 / MIT & Recomendado: Permite clasificar módulos críticos \\
        GPL v3 / AGPL & Restringido: Copyleft fuerte prohíbe cierre \\
        \bottomrule
    \end{tabular}
\end{table}

\begin{notaeditorial}
    \textbf{Nota Final}: El licenciamiento estratégico es tan importante como el desarrollo técnico para entornos de defensa y seguridad nacional.
\end{notaeditorial}

\end{document}
